\bookchapter{Глава 19: Диагностика неисправностей и самостоятельное спасение}

\chapterintro{Телефон, как и человек, иногда «болеет». Зависания, вылеты, чёрный экран, не заряжается... Большинство проблем не так страшны, как кажется, и вы вполне можете попробовать решить их самостоятельно, прежде чем нести в ремонт. Эта глава — ваш «справочник скорой помощи для телефона»: сначала пробуем недорогие методы, если не помогает — обращаемся к специалисту.}

\booksection{I. Телефон тормозит и греется — самый частый вопрос}

Честно говоря, то, что телефон после года использования начинает тормозить, — совершенно нормально. Причина не обязательно в поломке, чаще — в накопившемся «мусоре»: временные файлы, фоновые процессы, постепенно заполняющееся хранилище.

\booksubsection{Сначала диагностика, не спешите ремонтировать:}

\begin{enumerate}
  \item \textbf{Проверьте место в хранилище}: зайдите в Настройки Хранилище. Если осталось меньше 10\%, система будет сильно тормозить, пытаясь освободить место. Мой опыт: держите минимум 15\%–20\% свободного пространства для плавной работы.
  \item \textbf{Проверьте рейтинг расхода батареи}: зайдите в настройки батареи и посмотрите расход за последние 24 часа. Если какое-то приложение лидирует, но вы им почти не пользовались — значит, оно работает в фоне. Удалите или ограничьте фоновую активность.
  \item \textbf{Вспомните, что вы делали в последнее время}: тормоза начались после обновления системы? (Первые дни после крупного обновления торможение нормально — система перестраивает индексы в фоне.) Или после установки нового приложения? (Попробуйте удалить его.)
  \item \textbf{Оцените условия использования}: телефон на солнцепёке, в толстом чехле, с игрой на зарядке — перегрев неизбежен. Высокая температура заставляет процессор сбрасывать частоту — и вы чувствуете, что телефон внезапно стал медленным. Дайте остыть и попробуйте снова.
\end{enumerate}

\textbf{Методы лечения (от простых к сложным)}: перезагрузите телефон (решает 60\% проблем) освободите место (удалите ненужные приложения и видео) если виновато конкретное приложение, очистите его кэш или удалите и установите заново последнее средство: сброс до заводских настроек после резервного копирования.

\booksection{II. Приложения вылетают и не отвечают — когда приложение не слушается}

Приложение вылетает при запуске или зависает во время использования. Процесс похож на приём у врача — шаг за шагом:

\begin{enumerate}
  \item \textbf{Сначала закройте и откройте заново}: смахните приложение из фоновых и откройте снова. Большинство временных багов исчезает после этого.
  \item \textbf{Обновите приложение}: проверьте наличие обновлений в магазине приложений. Иногда проблема в том, что версия слишком старая и несовместима с системой.
  \item \textbf{Проверьте разрешения}: некоторым приложениям нужны определённые разрешения для работы, а вы их не дали. Например, приложение камеры без разрешения на камеру будет вылетать. Зайдите в настройки, найдите приложение и дайте нужные разрешения.
  \item \textbf{Очистите кэш}: Настройки Приложения найдите приложение Хранилище Очистить кэш. Внимание: «Очистить кэш», а не «Очистить данные»! Очистка данных удалит состояние входа и настройки приложения, фактически сбросив его до заводского состояния.
  \item \textbf{Удалить и переустановить}: последнее средство. Перед удалением убедитесь, что данные приложения не пропадут — для WeChat сначала сделайте резервную копию чатов, для игр убедитесь, что аккаунт привязан к номеру телефона.
\end{enumerate}

\booksection{III. Не заряжается или медленно заряжается — возможно, батарея не виновата}

Проблемы с зарядкой вызывают много тревоги, но часто причина в аксессуарах, а не в поломке телефона.

\begin{itemize}
  \item \textbf{Сначала проверьте порт зарядки}: аккуратно очистите порт зарядки от пыли и ворса зубочисткой или непроводящим инструментом. Я сам сталкивался: в кармане оторвалась нитка и забила порт зарядки. Казалось, что плохой контакт, а после очистки всё заработало.
  \item \textbf{Попробуйте другой кабель и зарядное устройство}: кабели выходят из строя гораздо чаще зарядных устройств, особенно те, что постоянно скручиваются. Проверьте с заведомо исправным кабелем и зарядным устройством — быстро определите, какой аксессуар неисправен.
  \item \textbf{Проблемы беспроводной зарядки}: если используете беспроводную зарядку, проверьте, не слишком ли толстый чехол и нет ли на нём металлических деталей (магнитные держатели и т.п.). И положите телефон точно по центру зарядной площадки — если смещён, может не заряжаться.
  \item \textbf{Медленная зарядка}: убедитесь, что мощность зарядного устройства достаточна (iPhone для быстрой зарядки требует минимум 20 Вт PD), кабель поддерживает быструю зарядку (не все кабели пропускают ток выше 3 А), телефон не перегрет (при перегреве мощность автоматически снижается для защиты батареи).
\end{itemize}

\booksection{IV. Сеть не подключается, Wi-Fi постоянно обрывается}

На улице мобильный интернет работает нормально, а дома Wi-Fi не подключается; или Wi-Fi подключён, но интернета нет. Проверяйте в таком порядке:

\begin{enumerate}
  \item \textbf{Переключите авиарежим}: откройте центр управления, включите авиарежим на 5 секунд и выключите. Это заставляет телефон заново искать сеть — очень эффективно.
  \item \textbf{«Забудьте» Wi-Fi и подключитесь заново}: в настройках Wi-Fi нажмите на значок «i» (на Android — шестерёнку) рядом с именем сети, выберите «Забыть эту сеть» и затем введите пароль заново.
  \item \textbf{Перезагрузите роутер}: у многих роутер не выключался месяцами, со временем это приводит к нестабильности. Отключите питание на 30 секунд и включите снова — большинство проблем с Wi-Fi решается.
  \item \textbf{Проверьте, в телефоне ли проблема}: другие устройства нормально подключаются к этому же Wi-Fi? Если да, а у вас не работает — возможно, телефону нужна перезагрузка или сброс настроек сети (Настройки Основные Перенос или сброс Сброс настроек сети — данные не удалятся, но все пароли Wi-Fi будут стёрты).
\end{enumerate}

\booksection{V. Проблемы с экраном: не реагирует на касания, мерцает, остаются следы}

\begin{itemize}
  \item \textbf{Плохая чувствительность касания}: сначала снимите защитную плёнку — возможно, дело в ней. Дешёвые стёкла со временем действительно ухудшают чувствительность. Вода на экране или жир на пальцах тоже вызывают дрейф касаний. Если нет — перезагрузите и попробуйте снова.
  \item \textbf{Следы/остаточное изображение на экране (выгорание)}: длительное отображение одной картинки (например, постоянно включённый экран с TikTok) может оставить «остаточное изображение» на OLED-экране. Это не поломка, а временное явление. Выключите экран на полчаса — обычно восстанавливается.
  \item \textbf{Мерцание экрана}: если мерцание при низкой яркости, возможно, проблема в ШИМ-регулировке, чаще встречается на OLED-экранах. Поищите в настройках «DC-регулировку» или «Защиту от мерцания» — включение этой опции заметно помогает. Не на всех телефонах есть эта функция.
  \item \textbf{Участок экрана полностью чёрный или белый}: это серьёзнее — скорее всего, физическое повреждение экрана, требуется замена. Экраны, заменённые в неофициальных сервисах, обычно уступают оригинальным по цветопередаче и яркости, а также нарушают водозащиту.
\end{itemize}

\booksection{VI. Не включается, завис на логотипе — не паникуйте}

Эта ситуация сильнее всего заставляет сердце биться чаще, но сохраняйте спокойствие и попробуйте следующее:

\begin{itemize}
  \item \textbf{Заряжайте полчаса и только потом пробуйте}: часто телефон просто полностью разрядился. Подключите оригинальное зарядное устройство, подождите минимум 30 минут и попробуйте включить. Некоторые телефоны после глубокого разряда требуют определённого уровня заряда для реакции.
  \item \textbf{Принудительная перезагрузка}: iPhone с Face ID — нажмите и отпустите кнопку увеличения громкости, затем кнопку уменьшения громкости, затем удерживайте боковую кнопку, пока не появится логотип Apple (может потребоваться 10–15 секунд). Android общий метод: удерживайте кнопку питания более 10 секунд до вибрации или включения экрана.
  \item \textbf{Циклическая перезагрузка на экране загрузки}: это называется «белое яблоко» (iPhone) или «зависание на логотипе» (Android). Если принудительная перезагрузка не помогает, возможно, потребуется режим восстановления через компьютер. iPhone: подключите к компьютеру с iTunes/Finder, войдите в режим восстановления и выберите «Обновить» (не «Восстановить») — система переустановится с сохранением данных. Android: войдите в режим Fastboot или Recovery, комбинации кнопок у каждого бренда свои — поищите «как войти в режим восстановления» для вашей модели.
  \item \textbf{Если ничего не помогло}: придётся обратиться к профессионалам. Если данные были сохранены, не стоит сильно переживать.
\end{itemize}

\booksection{VII. Проблемы со звуком: динамик не работает, тихий голос при разговоре}

\begin{itemize}
  \item \textbf{Динамик не издаёт звук}: сначала убедитесь, что это не беззвучный режим (переключатель на боковой грани iPhone), не режим «Не беспокоить» и не убавленная до нуля громкость. Звучит глупо, но такие случаи составляют 40\%. Затем проверьте, не подключён ли Bluetooth к каким-либо наушникам или колонке.
  \item \textbf{Собеседник вас не слышит}: проверьте, не забито ли отверстие микрофона — маленькие отверстия снизу телефона легко забиваются пылью. Аккуратно очистите мягкой щёткой. Не используйте иголки или зубочистки — можно повредить защитную сетку.
  \item \textbf{Хриплый/искажённый звук}: если в динамик попала вода — не сушите феном с горячим воздухом! Промокните влагу салфеткой и дайте высохнуть естественным путём. Приложения для выталкивания воды звуком (воспроизводят низкочастотный звук, чтобы вытолкнуть воду) имеют некоторый эффект — можно попробовать.
\end{itemize}

\booksection{VIII. Поиск официального ремонта: избегайте подводных камней}

Если все вышеперечисленные методы не помогли, нужен профессиональный ремонт. Но рынок ремонта полон pitfalls. Расскажу о собственном опыте и ошибках:

\begin{itemize}
  \item \textbf{Ищите настоящий официальный сервис, а не «похожий на официальный»}: зайдите на официальный сайт бренда (Apple.com.cn, Huawei.com.cn и т.д.) и найдите список «Авторизованных сервисных центров». Мастерские на улице с логотипом Apple, но без авторизации, могут использовать неоригинальные запчасти, и после ремонта водозащита исчезнет.
  \item \textbf{Что нужно сделать перед сдачей в ремонт}: сделайте резервную копию всех данных (не надейтесь, что данные не пропадут при ремонте), выйдите из всех важных аккаунтов (Apple ID, WeChat, Alipay), отключите «Локатор». Некоторые мастерские требуют пароль блокировки экрана — это риск для приватности. Лучше сами выйдите из аккаунтов перед передачей телефона.
  \item \textbf{Подводные камни при замене экрана}: неоригинальные экраны делятся на «перепрессованные» (оригинальное стекло заменено), «репрессованные» (неоригинальное стекло) и «сборные» (полностью неоригинальные). Цены отличаются в разы, но качество изображения и долговечность тоже сильно разнятся. Если телефон ещё на гарантии — идите официальным путём.
  \item \textbf{Подводные камни при замене батареи}: неоригинальные батареи могут иметь завышенную заявленную ёмкость и через полгода деградировать сильнее, чем оригинальная за два года. Замена батареи Apple в официальном сервисе ненамного дороже неофициального — рекомендую официальный путь.
  \item \textbf{Подводные камни при восстановлении данных}: если телефон намок или разбился, и в мастерской говорят «восстановление данных за дополнительную плату» — будьте осторожны. Недобросовестные мастера могут намеренно преувеличивать степень повреждения для завышения цены. Добросовестное восстановление данных: сначала оценка, потом работа, и плата только при успешном восстановлении.
\end{itemize}

\begin{tcbwarn}{Универсальный первый шаг: перезагрузка телефона}
С какой бы проблемой вы ни столкнулись — зависания, вылеты, проблемы с сетью, аномалии звука — сначала перезагрузите телефон. Это звучит как отговорка, но по моему опыту минимум 60\% мелких проблем исчезают после перезагрузки. Телефон, как и компьютер, со временем накапливает в памяти всякий мусор — перезагрузка очищает всё и начинает заново. Так что в следующий раз, когда телефон начнёт барахлить, не спешите искать «что делать» — перезагрузите и посмотрите.

\end{tcbwarn}

\begin{tcbpractice}{Запомните прямо сейчас: эти спасительные операции должны быть на автомате}
\begin{enumerate}
  \item \textbf{Принудительная перезагрузка}: iPhone с Face ID — нажать громкость вверх громкость вниз удерживать кнопку питания до появления логотипа. Android — удерживать кнопку питания 10–15 секунд до вибрации.
  \item \textbf{Очистка порта зарядки}: найдите маленький инструмент (зубочистку или скрепку для SIM-карты), аккуратно очистите порт зарядки от пыли.
  \item \textbf{Переподключение Wi-Fi}: в настройках Wi-Fi «забудьте» текущую сеть и введите пароль заново.
  \item \textbf{Включение/выключение авиарежима}: откройте центр управления, включите авиарежим, мысленно сосчитайте до пяти, выключите — проверьте, восстановилась ли сеть.
  \item Сделайте скриншоты этих шагов или запишите в заметки. Когда возникнет проблема, не нужно панически искать — просто откройте заметки.
\end{enumerate}

\end{tcbpractice}
